\begin{ex}
Tính giá trị của biểu thức: $125 + 45 - 70$.
\shortans{100}
\loigiai{
$170 - 70 = 100$.
}
\end{ex}

\begin{ex}
Tính giá trị của biểu thức: $80 - 20 + 35$.
\shortans{95}
\loigiai{
$60 + 35 = 95$.
}
\end{ex}

\begin{ex}
Tính giá trị của biểu thức: $24 \times 2 : 4$.
\shortans{12}
\loigiai{
$48 : 4 = 12$.
}
\end{ex}

\begin{ex}
Tính giá trị của biểu thức: $96 : 3 \times 2$.
\shortans{64}
\loigiai{
$32 \times 2 = 64$.
}
\end{ex}

\begin{ex}
Tính giá trị của biểu thức: $50 + 20 \times 3$.
\shortans{110}
\loigiai{
Nhân trước cộng sau: $50 + 60 = 110$.
}
\end{ex}

\begin{ex}
Tính giá trị của biểu thức: $90 - 40 : 2$.
\shortans{70}
\loigiai{
Chia trước trừ sau: $90 - 20 = 70$.
}
\end{ex}

\begin{ex}
Tính giá trị của biểu thức: $150 - 30 \times 4$.
\shortans{30}
\loigiai{
$150 - 120 = 30$.
}
\end{ex}

\begin{ex}
Tính giá trị của biểu thức: $200 + 400 : 5$.
\shortans{280}
\loigiai{
$200 + 80 = 280$.
}
\end{ex}

\begin{ex}
Tính giá trị của biểu thức: $(45 + 15) \times 3$.
\shortans{180}
\loigiai{
Trong ngoặc trước: $60 \times 3 = 180$.
}
\end{ex}

\begin{ex}
Tính giá trị của biểu thức: $(84 - 24) : 6$.
\shortans{10}
\loigiai{
Trong ngoặc trước: $60 : 6 = 10$.
}
\end{ex}

\begin{ex}
Tính giá trị của biểu thức: $120 : (2 + 4)$.
\shortans{20}
\loigiai{
$120 : 6 = 20$.
}
\end{ex}

\begin{ex}
Tính giá trị của biểu thức: $40 \times (10 - 6)$.
\shortans{160}
\loigiai{
$40 \times 4 = 160$.
}
\end{ex}

\begin{ex}
Tính giá trị của biểu thức: $35 + 25 \times 2 - 10$.
\shortans{75}
\loigiai{
$35 + 50 - 10 = 75$.
}
\end{ex}

\begin{ex}
Tính giá trị của biểu thức: $100 - 60 : 2 + 15$.
\shortans{85}
\loigiai{
$100 - 30 + 15 = 85$.
}
\end{ex}

\begin{ex}
Tính giá trị của biểu thức: $(30 + 50) : 4 + 12$.
\shortans{32}
\loigiai{
$80 : 4 + 12 = 20 + 12 = 32$.
}
\end{ex}

\begin{ex}
Tính giá trị của biểu thức: $250 - (60 + 40) \times 2$.
\shortans{50}
\loigiai{
$250 - 100 \times 2 = 250 - 200 = 50$.
}
\end{ex}

\begin{ex}
Tính giá trị của biểu thức: $48 : 4 + 36 : 3$.
\shortans{24}
\loigiai{
$12 + 12 = 24$.
}
\end{ex}

\begin{ex}
Tính giá trị của biểu thức: $15 \times 4 - 40 : 2$.
\shortans{40}
\loigiai{
$60 - 20 = 40$.
}
\end{ex}

\begin{ex}
Tính giá trị của biểu thức: $70 + 30 : (5 + 5)$.
\shortans{73}
\loigiai{
$70 + 30 : 10 = 70 + 3 = 73$.
}
\end{ex}

\begin{ex}
Tính giá trị của biểu thức: $(100 - 20) : (4 + 4)$.
\shortans{10}
\loigiai{
$80 : 8 = 10$.
}
\end{ex}

\begin{ex}
Tìm giá trị của $A$ biết $A = 12 \times 5 + 40$.
\shortans{100}
\loigiai{
$A = 60 + 40 = 100$.
}
\end{ex}

\begin{ex}
Tìm giá trị của $B$ biết $B = 180 : 3 - 25$.
\shortans{35}
\loigiai{
$B = 60 - 25 = 35$.
}
\end{ex}

\begin{ex}
Điền số thích hợp: $(35 + \dots) \times 2 = 100$.
\shortans{15}
\loigiai{
$35 + x = 50 \Rightarrow x = 15$.
}
\end{ex}

\begin{ex}
Điền số thích hợp: $80 - \dots : 3 = 60$.
\shortans{60}
\loigiai{
$x : 3 = 20 \Rightarrow x = 60$.
}
\end{ex}

\begin{ex}
Tính giá trị của biểu thức: $240 : 6 \times 5$.
\shortans{200}
\loigiai{
$40 \times 5 = 200$.
}
\end{ex}
